\documentclass[11pt]{article}
\usepackage[a4paper,margin=25mm]{geometry}
\usepackage[T1]{fontenc}
\usepackage{lmodern}
\usepackage{microtype}
\usepackage{amsmath,amssymb,amsthm,mathtools}
\usepackage{booktabs}
\usepackage{enumitem}
\usepackage{xcolor}
\usepackage[hidelinks]{hyperref}

\newtheorem{definition}{Definition}[section]
\newtheorem{theorem}[definition]{Theorem}
\newtheorem{lemma}[definition]{Lemma}
\newtheorem{proposition}[definition]{Proposition}
\newtheorem{corollary}[definition]{Corollary}
\newtheorem{assumption}[definition]{Assumption}
\newtheorem{remark}[definition]{Remark}

\newcommand{\Core}{\mathsf{Core}}
\newcommand{\BUOV}{\mathsf{BUOV}}
\newcommand{\Piiss}{\Pi_{\mathsf{issue}}}
\newcommand{\TEnc}{\mathsf{TEnc}}
\newcommand{\HDEM}{\mathsf{HDEM}}
\newcommand{\TPDec}{\mathsf{TPDec}}
\newcommand{\TComb}{\mathsf{TComb}}
\newcommand{\Verify}{\mathsf{Verify}}
\newcommand{\negl}{\mathsf{negl}}
\newcommand{\Adv}{\mathsf{Adv}}
\newcommand{\QPT}{\mathsf{QPT}}
\newcommand{\Enc}{\mathsf{Enc}}
\newcommand{\Dec}{\mathsf{Dec}}
\newcommand{\Hash}{\mathsf{H}}
\newcommand{\rid}{\mathsf{rid}}
\newcommand{\sn}{\mathsf{sn}}
\newcommand{\ctx}{\mathsf{ctx}}
\newcommand{\pp}{\mathsf{pp}}
\newcommand{\pk}{\mathsf{pk}}
\newcommand{\sk}{\mathsf{sk}}
\newcommand{\eps}{\varepsilon}

\title{A Post-Quantum Relation-Bound Blind Certification Core\\
with Signature-Gated Threshold Decryption}
\author{Research construction draft, version 1.8}
\date{23 August 2026}

\begin{document}
\maketitle

\begin{abstract}
This document specifies and proves an abstract post-quantum cryptographic core for
privacy-preserving, accountable credentials.  During an offline issuance phase, a
holder obtains a blind signature only after proving that the hidden ticket, the
blind-signing request, the holder secret, and a trace ciphertext are mutually
consistent.  During opening, threshold decryption shares are released only for a
validly signed ticket and an independently authorized case.  This
\emph{signature-gated threshold decryption} rule replaces a public validity proof
carried by every online ticket.

The main result is conditional and modular: assuming simulation-extractable
zero knowledge, Blind-UOV one-more unforgeability and
blindness, and a compact threshold encryption scheme with post-quantum
privacy and unique openings, the construction provides issuer unlinkability,
certified ciphertext validity, decoder safety, trace correctness,
non-frameability, and a gated chosen-ciphertext privacy notion against fewer than
the threshold number of opening authorities.  The result deliberately does not
claim ordinary IND-CCA security for the bare trace encryption.  Version 1.2
also instantiates the trace relation as a Niederreiter syndrome plus a
SHAKE256/KMAC256 one-time hash DEM, proves information-theoretic uniqueness of
the decoded error, and removes a redundant ticket commitment.  Version 1.3
freezes a paper-compatible Blind-UOV-Is/Shorter costing profile, fixes a
368-byte payload, gives a reproducible five-block circuit count, and separates
the large issuance NIZK from the CAP proof carried by the final Blind-UOV
signature.  Version 1.4 adds an executable witness builder, a streaming
characteristic-two circuit generator, independent SHAKE256/KMAC256 checks, and
full-circuit rejection tests.  Native Blind-UOV circuit integration, backend
simulation extractability, and the threshold implementation remain explicit
obligations.  Version 1.5 lowers the complete incremental relation into a
checksummed binary $\mathbb F_2$-R1CS archive, independently parses every row,
and separates the nonlinear-only cost model from the portable materialized-row
count.  Version 1.6 corrected the signer-view ABI but incorrectly treated two
independently openable 256-bit request lanes as a 512-bit random-function
output.  Version 1.7 removes that unsound amplification.  It selects one
Blind-UOV-III/Shorter instance with a 576-bit request target and derives
cross-message request binding in the QROM from CAP unique-witness soundness and
the collision bound for a shifted 576-bit random function.  The cached issuer
key grows to 189.2~KB, while the online ticket remains 12,012 bytes.  Native
Blind-UOV circuit integration and proof-backend qualification remain explicit
obligations.  Version 1.8 writes $h_{\BUOV}=H_{\BUOV}(m,c_r)$ and internalizes
the linear request equation $y=r+h_{\BUOV}$ as 576 binary equalities.  It
narrows the sole external boundary to
$h_{\BUOV}=H_{\BUOV}(m,\mathsf{CAP.Commit}(r;\rho))$ and adds a fail-closed
native-import contract.  It also freezes the paper's NIST-III/Shorter TCitH
tree parameters and exposes the required $\mathbb F_{2^{193}}$ backend lift.
\end{abstract}

\section{Status, scope, and non-goals}

This is the standalone cryptographic core.  It intentionally excludes satellite
AKE, radio transport, FLEO relay behavior, handover, revocation distribution, and
concrete parameter optimization.  Issuance and opening may use large proofs,
large public keys, and large private keys.  The constrained online path carries
only a ticket payload and one compact blind signature.

The construction is designed for the following division of trust.

\begin{itemize}[leftmargin=*]
  \item The issuing HNCC is \emph{honest but curious}: it authenticates the real
  identity, enforces a common issuance policy, verifies $\Piiss$, and signs only
  after acceptance, but later tries to link an observed ticket to an issuance
  session.
  \item Holders and network attackers may be malicious and quantum polynomial
  time.
  \item Fewer than $\tau$ of $n$ opening authorities may be corrupted.  Any
  authorized set of at least $\tau$ authorities can open a ticket.
  \item Issuer signing keys and opening decryption keys are independent even when
  the organizations operating them overlap.
  \item Per-user visible metadata is forbidden.  Version, epoch, policy,
  expiration bucket, and key identifiers are common to a sufficiently large
  anonymity set; otherwise they become an issuer watermark outside the
  cryptographic game.
\end{itemize}

\paragraph{Claim boundary.}
The theorems below are complete reductions for the abstract interfaces and games
defined here.  They are not a claim that the present Blind-UOV or threshold
Niederreiter papers already implement every required interface.  The Blind-UOV
request equations and their visibility are fixed in
Section~\ref{sec:buov-normal}; the exact added
circuit is counted in Section~\ref{sec:circuit-audit} and executed in
Section~\ref{sec:executable-relation}.  The resulting proof-size and prover
runtime figures remain paper-anchored estimates rather than implementation
benchmarks.  The executable generator is a relation validator, not a proof
backend.  The numerical Niederreiter set is a size anchor, not a deployment
approval.

\section{Notation and post-quantum conventions}

The security parameter is $\lambda$.  All adversaries and reductions are
uniform or non-uniform $\QPT$ algorithms.  A function $\negl(\lambda)$ is
negligible.  Encodings are canonical, length-prefixed, and domain separated.
Concatenation is written $\|$.  Every hash invocation includes a fixed protocol
label and version.

The issuing organization authenticates a real identity $\rid$ before the
cryptographic issuance protocol.  The session identifier $\mathsf{sid}$ is used
only for replay and quota control and is never copied into a ticket.  The common
issuance information is
\[
  \mathsf{ci}=(\mathsf{version},\mathsf{epoch},\mathsf{domain},
  \mathsf{policy},\mathsf{expiryBucket},\mathsf{kid}_{\mathrm{OA}},
  \mathsf{kid}_{\mathrm{I}}).
\]
Its fixed 32-byte ticket representation is
\[
 \ctx=\mathsf{SHAKE256}(\mathtt{PQ\mbox{-}RBBC/CTX}\|
       \mathsf{Encode}(\mathsf{ci}),256).
\]
The federation configuration signature authenticates $(\mathsf{ci},\ctx)$.
The digest is computed once outside $\Piiss$ and is a public input; the issuer
cannot select a per-user value.

For any game $G$, $\Adv^G_{\mathcal A}(\lambda)$ denotes the usual absolute
distance from the ideal winning probability.  Oracle calls are polynomially
bounded.

\section{Required primitive interfaces}

\subsection{Hashing and CAP commitment}

We use domain-separated hashes
\[
  H_{\mathrm{ticket}},\quad H_{\BUOV},\quad H_{\mathrm{hold}},\quad
  H_{\mathrm{evidence}}.
\]
They are modeled as quantum-resistant collision-resistant hashes.  A random
oracle or QROM instantiation may be used by a concrete component, but the core
proof calls only the explicitly named property.

For the version-1.3 circuit audit, fix
\begin{align*}
 H_{\mathrm{ticket}}(X)&=\mathsf{SHAKE256}
   (\mathtt{PQ\mbox{-}RBBC/TICKET}\|X,256),\\
 H_{\mathrm{hold}}(k)&=\mathsf{SHAKE256}
   (\mathtt{PQ\mbox{-}RBBC/HOLD}\|k,256).
\end{align*}
The Blind-UOV hash $H_{\BUOV}$ and CAP commitment internals use the paper's
Anemoi-based costing profile; they are already included in its baseline NIZK
count.  This mixed choice retains standardized SHAKE/KMAC for new protocol
logic while avoiding the multi-billion-constraint SHA3+AES baseline inside the
Blind-UOV commitment proof.

There is no separate public commitment to the ticket payload.  The accepting
proof extracts $M$, Equation~\eqref{eq:i2} fixes its digest, and the Blind-UOV
request binds that digest.  Publishing an additional $c_M$ would add a primitive
and a circuit block without strengthening this chain.

\subsection{Blind-UOV request normal form and visibility}\label{sec:buov-normal}

The blind-signature interface is
\[
  (\mathsf{bsk},\mathsf{bvk})\leftarrow\BUOV.\mathsf{KeyGen}(1^\lambda),
\]
with an interactive request/sign/unblind protocol and a public verification
algorithm
\[
  \BUOV.\mathsf{Vfy}(\mathsf{bvk},m,\sigma).
\]

The core needs more than ordinary blindness: the issuance proof must identify
the exact hidden message represented by the request.  Protocols 3 and 4 of the
31 October 2025 revision have the user compute~\cite{blind-uov}
\begin{align}
 c_r&\leftarrow \mathsf{CAP.Commit}(r;\rho),
   \tag{B1}\label{eq:b1}\\
 y&=r+H_{\BUOV}(m,c_r)\in\mathbb F_q^o.
   \tag{B2}\label{eq:b2}
\end{align}
Only $(y,\pi_1)$ is sent to the signer.  The values $m,r,\rho,c_r$ are
witness/state, not signer-visible request fields.  The final signature later
contains $(c_r,c_x,\pi_2)$.  Publishing $c_r$ during issuance would both depart
from the paper and give the issuer a direct link to the final signature.

A Blind-UOV-Is target is only 256 bits.  Our authorization-binding property is
a freely chosen cross-message collision problem, not a fixed-target preimage
problem, so its generic quantum scale is only about
$2^{256/3}=2^{85}$~\cite{bht,compressed-oracle}.  Version 1.7 instead selects
the paper's Blind-UOV-III/Shorter profile: $r$ and $y$ are 576 bits, the cached
public key is 189.2~KB, and the final signature is 11,644 bytes
\cite{blind-uov}.  Define
\[
 J_m(r,\rho)=r+H_{\BUOV}
   (m,\mathsf{CAP.Commit}(r;\rho)).
\]
The public request is exactly the 72-byte value $\beta=y$; its hidden witness
is $b=(r,\rho)$.

\begin{definition}[Blind-UOV-III request relation]\label{def:rrbs}
For $\beta=y$ and $b=(r,\rho)$, define
\[
 R^{\BUOV\text{-}III}_{\mathsf{req}}(y;m,r,\rho)=1
 \Longleftrightarrow
 c_r=\mathsf{CAP.Commit}(r;\rho)\ \wedge\
 y=r+H_{\BUOV}(m,c_r),
\]
where $c_r$ is a derived hidden witness.  A completed session returns one
ordinary Blind-UOV-III signature $\sigma=(c_r,c_x,\pi_2)$ on $m$.
\end{definition}

\begin{theorem}[QROM cross-message request binding]\label{thm:j-qrom}
Let $H_{\BUOV}:\{0,1\}^*\rightarrow\{0,1\}^{576}$ be an independent quantum
random oracle.  Assume the CAP commitment has straightline extraction and a
unique committed mask: producing one admissible $c_r$ with two different
extracted $r$ values is negligible.  If an adversary makes at most $q$ relevant
quantum-oracle/evaluation queries, then
\[
 \Adv^{\mathsf{xmsg}}_{J}(\mathcal A)
 \leq
 \Adv^{\mathsf{uw}}_{\mathsf{CAP}}(\mathcal B_1)
 +\Adv^{\mathsf{ext}}_{\mathsf{CAP}}(\mathcal B_2)
 +O\!\left(\frac{q^3}{2^{576}}\right).
\]
In particular, the generic quantum collision scale is $2^{576/3}=2^{192}$;
at $q=2^{128}$ the random-function term is $O(2^{-192})$.
\end{theorem}

\begin{proof}
Run the straightline CAP extractor on every admissible commitment relevant to
the adversary's output.  Abort on extraction failure.  If one commitment yields
two different masks, output a CAP unique-witness violation.  Otherwise each
admissible commitment $c$ has a well-defined value $r(c)$ fixed before the
domain-separated query $H_{\BUOV}(m,c)$.  Quotienting duplicate openings with
the same $(c,r(c))$, define
\[
 Z(m,c)=H_{\BUOV}(m,c)+r(c).
\]
For every admissible input, addition of the fixed $r(c)$ is a permutation of
the 576-bit oracle range.  Thus $Z$ is a pointwise-shifted 576-bit random
function.  If the adversary gives two accepted openings of one public $y$ for
$m\neq m'$, then $(m,c)\neq(m',c')$ and
$Z(m,c)=Z(m',c')=y$, a collision.  The compressed-oracle collision bound is
$O(q^3/2^{576})$~\cite{compressed-oracle}; the stated CAP bad events cover the
remaining cases.
\end{proof}

\begin{remark}[Why v1.6 parallel repetition was removed]
Two independently openable 256-bit lanes do not behave as one monolithic
512-bit random function.  After fixing $m,m'$, an attacker can run a BHT claw
search for lane 0 and then another for lane 1, for total work about
$2\cdot2^{85}$ rather than $2^{170}$.  A larger single request range gives the
required collision exponent without adding a correlation assumption that might
damage blindness.
\end{remark}

\begin{remark}[Why ordinary blind signing is insufficient]
If $\Piiss$ proves a valid ciphertext but the proof is not tied to the exact
blind request $\beta$, a malicious holder can prove statement $M_0$ and ask for
a signature on $M_1$.  The
$R^{\BUOV\text{-}III}_{\mathsf{req}}$ conjunct is the formal
connection that prevents this substitution.
\end{remark}

\subsection{Simulation-extractable NIZK argument of knowledge}

The proof system $\Piiss=(\mathsf{Setup},\mathsf{Prove},\mathsf{Vfy})$ is a
post-quantum simulation-extractable NIZK argument of knowledge for the issuance
relation in Section~\ref{sec:issue-relation}.  Its setup produces a CRS; the
security proof may use simulation and extraction trapdoors.  In the real
protocol no trapdoor is available to any party.

The proof may be several megabytes because it is transmitted only during
infrequent issuance.  This design choice is intentional.  Knowledge
soundness alone is enough for decoder safety in a simple stand-alone execution;
simulation extractability is used to support concurrent sessions and the
gated-CCA hybrid proof.

\subsection{Compact threshold trace encryption}

The threshold encryption interface is
\begin{align*}
  (\mathsf{tpk},\mathsf{tsk}_1,\ldots,\mathsf{tsk}_n)
    &\leftarrow \TEnc.\mathsf{KeyGen}(1^\lambda,n,\tau),\\
  C&\leftarrow \TEnc.\mathsf{Enc}(\mathsf{tpk},\mu;\mathsf{ad},\rho),\\
  d_i&\leftarrow \TPDec(\mathsf{tsk}_i,C),\\
  \mu\text{ or }\bot&\leftarrow \TComb(C,\{d_i\}_{i\in I}),\qquad |I|\geq\tau.
\end{align*}

There is an NP relation
\[
 R_{\mathsf{enc}}(\mathsf{tpk},C,\mathsf{ad};\mu,\rho)=1
 \quad\Longleftrightarrow\quad
 C=\TEnc.\mathsf{Enc}(\mathsf{tpk},\mu;\mathsf{ad},\rho).
\]
Only ciphertexts in the language $L_{\mathsf{enc}}$ are promised to be safe
inputs to the distributed decoder.

We require the following post-quantum properties.

\begin{enumerate}[label=(T\arabic*),leftmargin=*]
  \item \textbf{Correctness on valid inputs.}  Combining any $\tau$ correct
  shares for a ciphertext in $L_{\mathsf{enc}}$ returns its plaintext.
  \item \textbf{$\tau$-privacy under chosen plaintext.}  A $\QPT$ adversary
  holding fewer than $\tau$ secret shares cannot distinguish encryptions of two
  equal-length plaintexts.
  \item \textbf{Unique opening.}  It is infeasible to find
  $C,\mu,\rho,\mu',\rho'$ with $\mu\neq\mu'$ such that both
  $R_{\mathsf{enc}}(\mathsf{tpk},C,\mathsf{ad};\mu,\rho)=1$ and
  $R_{\mathsf{enc}}(\mathsf{tpk},C,\mathsf{ad};\mu',\rho')=1$.
  \item \textbf{Robust shares.}  Invalid or inconsistent partial decryptions are
  detected by $\TComb$ or by a component proof/MAC before a plaintext is
  accepted.
\end{enumerate}

Ordinary IND-CCA security of the bare encryption is not required and is not
claimed.

\subsection{Opening authorization}

An independent post-quantum authorization mechanism verifies
\[
  \mathsf{GateAuthVfy}(\mathsf{gavk},d_M,\mathsf{caseID},
  H_{\mathrm{evidence}}(E),\mathsf{auth})=1,
\]
where $d_M=H_{\mathrm{ticket}}(\mathsf{Encode}(M))$.  The authorization binds a
ticket digest, case identifier, evidence digest, expiry, and purpose.  Its
unforgeability and policy governance are orthogonal to trace-ciphertext privacy,
but the binding is required to prevent using an authorization for a different
ticket.

\section{Ticket format and issuance relation}\label{sec:issue-relation}

The hidden ticket payload is
\[
 M=(\ctx,\sn,h,C),
 \qquad h=H_{\mathrm{hold}}(k_{\mathrm{hold}}),
 \qquad C=\mathsf{Ntr.Enc}(\mathsf{tpk},\ctx,h;\rid,\sn,e).
\]
For the v1.3 profile all fields are fixed width and
\[
 \mathsf{Encode}(M)=\ctx\|\sn\|h\|s\|c_\mu\|\gamma
 \in\{0,1\}^{2944},
\]
namely $32+16+32+208+48+32=368$ bytes.  No variable-length field or
user-selected policy string occurs in the payload.
The serial number $\sn\leftarrow\{0,1\}^{128}$ and holder key
$k_{\mathrm{hold}}\leftarrow\{0,1\}^{256}$ are fresh.  The serial appears both
outside and inside the trace ciphertext so that an opener can detect payload
splicing.  The holder secret itself never appears in the ticket.  Define
\[
  m=H_{\mathrm{ticket}}(\mathsf{Encode}(M)).
\]

The public issuance statement and private witness are
\begin{align*}
x&=(\pp,\ctx,\mathsf{sid},\rid,\beta),\\
w&=(M,r,\rho,k_{\mathrm{hold}},e).
\end{align*}
The relation $R_{\mathsf{issue}}(x;w)=1$ iff all of the following hold:
\begin{align}
 \mathsf{TicketShape}(M,\ctx)&=1, \tag{I1}\label{eq:i1}\\
 m&=H_{\mathrm{ticket}}(\mathsf{Encode}(M)), \tag{I2}\label{eq:i2}\\
 R^{\BUOV\text{-}III}_{\mathsf{req}}(\beta;m,r,\rho)&=1,
   \tag{I3}\label{eq:i3}\\
 M.h&=H_{\mathrm{hold}}(k_{\mathrm{hold}}), \tag{I4}\label{eq:i4}\\
 R_{\mathsf{enc}}(\mathsf{tpk},M.C,
   \mathsf{Encode}(\ctx,M.\sn,M.h);\rid\|M.\sn,e)&=1.
   \tag{I5}\label{eq:i5}
\end{align}
The real identity and session identifier are public to the issuer during
issuance but are not part of $M$.  Equation~\eqref{eq:i5} binds the hidden
trace plaintext to the identity authenticated in that exact session.

\section{Construction}

\subsection{Setup}

$\Core.\mathsf{Setup}(1^\lambda,n,\tau)$ performs:
\begin{enumerate}[leftmargin=*]
  \item generate $\Piiss$ parameters and the domain-separated hash parameters;
  \item generate issuer keys
  $(\mathsf{bsk},\mathsf{bvk})\leftarrow\BUOV.\mathsf{KeyGen}(1^\lambda)$;
  \item generate independent opening keys
  $(\mathsf{tpk},\mathsf{tsk}_1,\ldots,\mathsf{tsk}_n)
  \leftarrow\TEnc.\mathsf{KeyGen}(1^\lambda,n,\tau)$;
  \item publish $(\mathsf{ci},\ctx)$ authenticated by a federation
  configuration signature.
\end{enumerate}

No party is allowed to derive $\mathsf{bsk}$ and any $\mathsf{tsk}_i$ from the
same seed or key hierarchy.

\subsection{Blind relation-bound issuance}

For authenticated identity $\rid$ and fresh issuer-side $\mathsf{sid}$:

\begin{enumerate}[leftmargin=*]
  \item The holder samples $\sn,k_{\mathrm{hold}},e$ and computes
  \[
   h=H_{\mathrm{hold}}(k_{\mathrm{hold}}),\qquad
   C=\mathsf{Ntr.Enc}(\mathsf{tpk},\ctx,h;\rid,\sn,e)
  \]
  according to Equations~\eqref{eq:n1}--\eqref{eq:n7},
  constructs $M$, and
  computes $m=H_{\mathrm{ticket}}(\mathsf{Encode}(M))$.
  \item It samples $r\in\{0,1\}^{576}$ and $\rho$, computes
  \[
    c_r=\mathsf{CAP.Commit}(r;\rho),\qquad
    y=r+H_{\BUOV}(m,c_r),
  \]
  sets the public request $\beta=y$, and produces
  \[
     \pi_{\mathsf{issue}}
     \leftarrow \Piiss.\mathsf{Prove}(\mathsf{crs},x,w).
  \]
  \item The issuer checks authenticated identity, epoch, quota, replay status,
  and
  \[
    \Piiss.\mathsf{Vfy}(\mathsf{crs},x,\pi_{\mathsf{issue}})=1.
  \]
  Only then does it execute
  $z\leftarrow\BUOV.\mathsf{Respond}(\mathsf{bsk},y)$.
  \item The holder computes
  \[
    \sigma\leftarrow
    \BUOV.\mathsf{Finalize}(m,c_r,y,z;r,\rho)
  \]
  and accepts only if
  $\BUOV.\mathsf{Vfy}(\mathsf{bvk},m,\sigma)=1$.
\end{enumerate}

The proof $\pi_{\mathsf{issue}}$ is not an extra proof placed beside the
Blind-UOV well-built-request proof.  It is the strengthened replacement for
that proof: the original request equations are retained and the ticket,
identity, holder, and trace-ciphertext conjuncts are appended to the proved
relation.  Thus issuance sends one augmented relation proof.

The spendable ticket is $T=(M,\sigma)$.  Neither $\beta$,
$\pi_{\mathsf{issue}}$, $\rid$, nor $\mathsf{sid}$ is transmitted with $T$.

\subsection{Ticket verification}

$\Core.\mathsf{VerifyTicket}(T)$ recomputes
$m=H_{\mathrm{ticket}}(\mathsf{Encode}(M))$, checks canonical payload
encoding and the public common-information policy, then returns
$\BUOV.\mathsf{Vfy}(\mathsf{bvk},m,\sigma)$.

This check proves neither public ciphertext validity nor identity correctness by
itself.  Those facts follow from the certified-issuance theorem and the trust
assumption that the issuer signs only after accepting $\Piiss$.

\subsection{Signature-gated threshold opening}

For an opening request
\[
  Q=(T,E,\mathsf{caseID},\mathsf{auth}),
\]
each honest opening authority runs the following gate before touching its
decoder:

\begin{enumerate}[leftmargin=*]
  \item require $\Core.\mathsf{VerifyTicket}(T)=1$;
  \item require the common epoch and key identifiers to be acceptable;
  \item compute $d_M=H_{\mathrm{ticket}}(\mathsf{Encode}(M))$ and require a
  valid, unused, purpose-limited authorization bound to
  $(d_M,\mathsf{caseID},H_{\mathrm{evidence}}(E))$;
  \item only after all checks pass, output
  $d_i=\TPDec(\mathsf{tsk}_i,M.C)$, bound to the same request digest.
\end{enumerate}

The combiner accepts a reconstructed plaintext $\rid'\|\sn'$ only if robust
share verification succeeds and $\sn'=M.\sn$.  Otherwise it outputs $\bot$.

\paragraph{Critical API rule.}
The implementation MUST NOT export
$\TPDec(\mathsf{tsk}_i,C)$ as a callable production endpoint.  The only exposed
operation is $\Core.\mathsf{OpenShare}(\mathsf{tsk}_i,Q)$, which internally
performs the gate.  This is a security boundary, not an application convention.

\section{Security games}

\subsection{Certified-language soundness}

The challenger maintains a log $\mathcal L$ of accepted issuance sessions.  In
the analysis it uses the extractor to associate each completed session with
$(\rid,M,w)$.  An adversary wins $\mathsf{CertSound}$ if it outputs a ticket
$T=(M,\sigma)$ for which $\Core.\mathsf{VerifyTicket}(T)=1$ and either:
\begin{enumerate}[label=(\alph*),leftmargin=*]
  \item $M$ is not associated with any completed accepted issuance in
  $\mathcal L$; or
  \item every associated extraction fails one of
  Equations~\eqref{eq:i1}--\eqref{eq:i5}.
\end{enumerate}

Distinct messages are counted by their canonical ticket digest.  This game
captures both a new signature and a proof/request substitution.

\subsection{Issuer unlinkability}

An honest-protocol curious issuer chooses or observes two authenticated
identities $\rid_0,\rid_1$ that receive the same $\mathsf{ci}$.  Two honest
holders execute issuance concurrently.  The issuer obtains its full protocol
views, after which it receives the two resulting tickets in random order and
guesses the order.  It may retain all issuance state but receives no opening
shares.  The advantage over $1/2$ is issuer-unlinkability advantage.

This game allows the issuer to know which identity participated in each
issuance; it hides which later ticket came from which session.  It cannot hide
timing, a uniquely sized policy, or a per-user expiration bucket, which is why
common metadata is a stated system requirement.

\subsection{Gated chosen-ciphertext trace privacy}\label{sec:gcca-game}

The game $\mathsf{t\text{-}GCCA}$ is parameterized by threshold $\tau$.
The adversary receives all public parameters and fewer than $\tau$ opening
secret shares.  It may run malicious-user issuance sessions for identities it
is authorized to use, receiving completed tickets, and may submit opening
queries for any validly authorized ticket other than the exact challenge ticket
digest.

For the challenge it selects equal-length identities $\rid_0,\rid_1$ and common
ticket metadata.  The challenger samples $b\leftarrow\{0,1\}$, forms a valid
ticket whose trace ciphertext encrypts $\rid_b\|\sn^*$, and returns
$T^*=(M^*,\sigma^*)$.  The adversary wins by guessing $b$.

The HNCC issuer is not corrupted in this game because it necessarily observes
the authenticated identity at issuance.  HNCC privacy is instead covered by
issuer unlinkability.  The game models an outside observer, an evidence
processor, and fewer than $\tau$ opening authorities.

\begin{remark}[Why the game is not ordinary CCA]
The opening oracle accepts only issuer-certified ticket payloads, not arbitrary
ciphertexts.  Calling the result IND-CCA for bare $\TEnc$ would be incorrect.
The security statement is protocol-specific: certification plus the gate turns
chosen-ciphertext access into chosen-\emph{certified}-ciphertext access.
\end{remark}

\subsection{Trace soundness and non-frameability}

An adversary wins $\mathsf{TraceSound}$ if a valid ticket produced from a
completed accepted issuance authenticated as $\rid$ opens to
$\rid'\neq\rid$, or if it opens with an external serial number different from
the encrypted serial.  The issuer is honest in following the protocol.  The
holder may be malicious.

\section{Security results}

\subsection{Correctness}

\begin{theorem}[End-to-end correctness]\label{thm:correct}
For an honest issuance authenticated as $\rid$, an honestly produced ticket
$T=(M,\sigma)$ verifies.  Any authorized set of at least $\tau$ honest opening
authorities reconstructs $\rid\|M.\sn$, except with the sum of the component
correctness errors.
\end{theorem}

\begin{proof}
Honest computation satisfies every conjunct of $R_{\mathsf{issue}}$, so
$\Piiss$ accepts.  Blind-signature correctness gives a verifying signature
on $H_{\mathrm{ticket}}(\mathsf{Encode}(M))$.
Equation~\eqref{eq:i5} gives a
valid encryption of $\rid\|M.\sn$.  Threshold-encryption correctness on valid
inputs reconstructs that plaintext from any $\tau$ correct shares.  The serial
cross-check therefore succeeds.
\end{proof}

\subsection{Blind-signing message binding}

\begin{lemma}[Request uniqueness]\label{lem:buov-unique}
Let $\beta=y$.  If the Blind-UOV-III request relation accepts valid openings
for both $m$ and $m'\neq m$, then those values constitute a winning output in
the cross-message game of Theorem~\ref{thm:j-qrom}.
\end{lemma}

\begin{proof}
Expanding the two accepted relations gives distinct inputs $(m,c_r)$ and
$(m',c'_r)$ whose shifted-oracle values are both $y$.  This is precisely the
cross-message collision covered by Theorem~\ref{thm:j-qrom}, including the CAP
extraction and unique-witness bad events.
\end{proof}

\begin{lemma}[Request/proof binding]\label{lem:reqbind}
Suppose an issuance proof for $\beta$ is accepting and extractable.  Except
with cross-message request-collision probability, the message signed through
$\beta$ is exactly $H_{\mathrm{ticket}}(\mathsf{Encode}(M))$ for the payload
$M$ in the extracted witness.  A different payload cannot verify under the same
signed digest except by finding a collision in $H_{\mathrm{ticket}}$.
\end{lemma}

\begin{proof}
Extraction yields $(M,r,\rho,\ldots)$ satisfying
Equations~\eqref{eq:i2}--\eqref{eq:i3}.  Equation~\eqref{eq:i2} fixes $m$ to the
canonical ticket digest.  Lemma~\ref{lem:buov-unique} makes the message
represented by $\beta$ equal to this $m$.  If a different canonical payload
$M'\neq M$ verifies under that signed value, then
$H_{\mathrm{ticket}}(\mathsf{Encode}(M'))=
H_{\mathrm{ticket}}(\mathsf{Encode}(M))$, a collision.
\end{proof}

\begin{theorem}[Certified-language soundness]\label{thm:certsound}
For any $\QPT$ adversary $\mathcal A$, there exist $\QPT$ reductions
$\mathcal B_1,\ldots,\mathcal B_4$ such that
\begin{align*}
 \Adv^{\mathsf{CertSound}}_{\Core,\mathcal A}
 \leq{}&
 \Adv^{\mathsf{SE\text{-}AoK}}_{\Piiss,\mathcal B_1}
 +\Adv^{\mathsf{1mUF}}_{\BUOV,\mathcal B_2}
 +\Adv^{\mathsf{xmsg}}_{J,\mathcal B_3}
 +\Adv^{\mathsf{coll}}_{H_{\mathrm{ticket}},\mathcal B_4}.
\end{align*}
By Theorem~\ref{thm:j-qrom}, the third term is at most the two CAP bad-event
advantages plus $O(q^3/2^{576})$.
Consequently, every verifying ticket is associated with an accepted issuance
witness satisfying $R_{\mathsf{issue}}$, except with negligible probability.
\end{theorem}

\begin{proof}
Consider a verifying output $(M,\sigma)$.  If its canonical digest was not the
message of a completed blind-signing session, then the output is a one-more
forgery.  Otherwise consider the corresponding accepting issuance proof.
Simulation extraction returns a witness or violates the SE-AoK property.
Lemma~\ref{lem:reqbind} associates the signed digest with the extracted payload,
except under the remaining claw and collision bad events.  The extracted witness
then satisfies every issuance conjunct, including ciphertext validity and
identity binding.
\end{proof}

\begin{corollary}[Certified ciphertext validity]\label{cor:validct}
If $\Core.\mathsf{VerifyTicket}(M,\sigma)=1$, then $M.C\in
L_{\mathsf{enc}}$ and it has an opening to the authenticated issuance identity
and the external serial number, except with the bound in
Theorem~\ref{thm:certsound}.
\end{corollary}

\subsection{Decoder safety}

\begin{theorem}[Signature-gated decoder safety]\label{thm:safety}
Assume honest opening authorities expose only
$\Core.\mathsf{OpenShare}$ and execute its checks in the specified order.  The
probability that an honest authority evaluates $\TPDec$ on a ciphertext outside
$L_{\mathsf{enc}}$ is at most
\[
 \Adv^{\mathsf{CertSound}}_{\Core}+\Adv^{\mathsf{forge}}_{\mathsf{GateAuth}}
 +\Adv^{\mathsf{coll}}_{H_{\mathrm{ticket}}}.
\]
\end{theorem}

\begin{proof}
Reaching $\TPDec$ implies that the ticket signature and the opening
authorization both verified for the same ticket digest.  By
Corollary~\ref{cor:validct}, a verifying ticket contains a ciphertext in
$L_{\mathsf{enc}}$.  A mismatch between the ticket checked by the gate and the
ciphertext passed to the decoder either violates the atomic API invariant or
requires a digest collision.  A request not authorized for that digest is an
authorization forgery.  Taking a union bound gives the result.
\end{proof}

\begin{remark}
This theorem explains what is gained by moving the validity proof out of the
online ticket.  The large issuance proof is amortized, while the small issuer
signature becomes a certificate that the ciphertext was once proved valid.
The theorem fails immediately if an implementation offers a raw partial-
decryption endpoint.
\end{remark}

\subsection{Issuer unlinkability}

\begin{theorem}[Issuer unlinkability]\label{thm:unlink}
For an honest-protocol curious issuer, common visible metadata, and two honest
holders, there exist $\QPT$ reductions such that
\begin{align*}
 \Adv^{\mathsf{unlink}}_{\Core}
 \leq{}& \Adv^{\mathsf{blind}}_{\BUOV}
 +2\Adv^{\mathsf{zk}}_{\Piiss}
 +\Adv^{\mathsf{t\text{-}IND\text{-}CPA}}_{\TEnc}
 +\Adv^{\mathsf{coll}}_{H_{\mathrm{ticket}}}.
\end{align*}
\end{theorem}

\begin{proof}
First replace each real issuance proof with a simulated proof; zero knowledge
bounds the change.  The remaining request/response views are ordinary
Blind-UOV views, so signer blindness permits swapping the two hidden messages
without the issuer detecting the swap.  Threshold-encryption privacy prevents the issuer, which
holds no opening threshold, from matching a final trace ciphertext to either
known issuance identity.  Canonical digest collisions are excluded.  After
these replacements, the issuer's view is independent of the final permutation,
so its success
probability is $1/2$.
\end{proof}

\paragraph{Operational qualifier.}
The theorem is cryptographic, not traffic-analytic.  Issuance batching, fixed
message sizes, delayed activation, and common metadata are required to keep
timing and formatting from becoming a link tag.

\subsection{Trace correctness and non-frameability}

\begin{theorem}[Trace soundness]\label{thm:trace}
Under certified-language soundness and unique opening of $\TEnc$, a malicious
holder completing issuance while authenticated as $\rid$ cannot produce a
verifying ticket that opens to $\rid'\neq\rid$ or to a different serial number,
except with negligible probability.
\end{theorem}

\begin{proof}
By Theorem~\ref{thm:certsound}, a verifying ticket is associated with an
extracted accepted issuance witness.  Equation~\eqref{eq:i5} states that its
ciphertext opens to $\rid\|M.\sn$.  If threshold opening returns a different
plaintext while all robust-share checks pass, either threshold correctness
fails or the ciphertext has two different valid openings.  The latter violates
unique opening.  The explicit serial comparison catches any external/internal
serial mismatch.
\end{proof}

\begin{remark}[Issuer framing]
The theorem assumes an honest-protocol issuer because the issuer learns $\rid$
and controls the signing key.  A malicious issuer could directly sign a payload
encrypting an arbitrary victim identity.  Removing that trust requires a
threshold blind issuer, an auditable issuance log with privacy-preserving
accountability, or retaining a publicly checked proof.  It is not solved by the
present compact-ticket core.
\end{remark}

\subsection{Gated chosen-ciphertext trace privacy}

\begin{theorem}[$\tau$-GCCA trace privacy]\label{thm:gcca}
Assume fewer than $\tau$ opening authorities are corrupted.  For every $\QPT$
adversary $\mathcal A$ in the game of Section~\ref{sec:gcca-game}, there exist $\QPT$
reductions $\mathcal B_i$ such that
\begin{align*}
 \Adv^{\mathsf{t\text{-}GCCA}}_{\Core,\mathcal A}
 \leq{}&
 \Adv^{\mathsf{t\text{-}IND\text{-}CPA}}_{\TEnc,\mathcal B_0}
 +\Adv^{\mathsf{SE\text{-}AoK}}_{\Piiss,\mathcal B_1}
 +\Adv^{\mathsf{1mUF}}_{\BUOV,\mathcal B_2}\\
 &+\Adv^{\mathsf{unique}}_{\TEnc,\mathcal B_3}
 +\Adv^{\mathsf{xmsg}}_{J,\mathcal B_4}\\
 &
 +\Adv^{\mathsf{coll}}_{H_{\mathrm{ticket}},\mathcal B_5}.
\end{align*}
\end{theorem}

\begin{proof}
We use a sequence of games.

\paragraph{Game 0.}
This is the real $\mathsf{t\text{-}GCCA}$ game.

\paragraph{Game 1: extract accepted issuance sessions.}
For every accepted adversarial issuance proof, use the SE-AoK extractor to
record
\[
  (d_M,M,\rid,e,k_{\mathrm{hold}},r,\rho,\ldots)
\]
in a private table $\mathcal L$.  Abort if an accepting proof has no valid
witness.  The difference from Game 0 is bounded by SE-AoK security.  Request
binding and hash collision resistance ensure that the final signed digest is
the digest recorded in $\mathcal L$.

\paragraph{Game 2: reject unlogged verifying tickets.}
If the adversary submits a verifying ticket whose digest is neither the
challenge digest nor in $\mathcal L$, abort.  Such a ticket supplies an
additional valid message/signature pair not obtained from a completed issuance,
which yields a one-more forgery.  Proof/request substitution and digest
ambiguity are charged to the corresponding terms in the theorem.

\paragraph{Game 3: simulate the opening oracle from the issuance table.}
For every non-challenge verifying ticket, look up the extracted plaintext in
$\mathcal L$ and return it after the authorization checks.  This has the same
result as threshold opening by correctness and unique opening, but it uses no
honest opening secret share and never invokes a decryption oracle.  A ticket
that reuses the challenge ciphertext under a newly certified digest requires an
accepted proof containing a valid opening of that ciphertext.  Extraction then
reveals the challenge plaintext.  If the extracted opening disagrees with the
challenge plaintext, it violates unique opening; if it agrees, the reduction
can immediately guess the threshold-encryption challenge bit.  Thus challenge
re-certification does not create an unanswerable decryption query.

\paragraph{Game 4: replace the challenge encryption.}
Embed the $\tau$-IND-CPA challenge of $\TEnc$ as $M^*.C$.  The reduction is
given fewer than $\tau$ shares exactly as in the base game.  It samples the
remaining visible ticket fields identically.  Because it generated the issuer
key and knows the complete challenge message $m^*$, it locally executes the
genuine Blind-UOV protocol and obtains one exactly distributed final signature;
it does not need the unknown trace-encryption randomness.
The challenge issuance transcript is not given to the adversary in this game,
and $\pi_{\mathsf{issue}}$ influences only the accept/reject gate, not the
distribution of the response after acceptance.  Therefore the reduction may
internally bypass that gate when manufacturing the challenge ticket without
changing any value visible to the adversary.
All opening queries are answered from $\mathcal L$ as in Game 3.  Therefore
changing the encrypted
identity from $\rid_0$ to $\rid_1$ changes the adversary's view by at most the
$\tau$-IND-CPA advantage of $\TEnc$.

In the final game the challenge bit is independent of the adversary's view.
Summing the hybrid differences proves the theorem.
\end{proof}

\paragraph{Why the reduction is useful.}
The large issuance proof is not merely a validity check.  Its extractor gives
the simulator the plaintext and randomness of every non-challenge ciphertext
that can ever obtain a valid issuer signature.  Therefore the simulator can
answer gated opening queries without a base decryption oracle.  This is the
precise reason a CPA-secure compact threshold encryption can suffice inside the
gated protocol even though it would not suffice under an unrestricted CCA API.

\subsection{One-more ticket unforgeability}

\begin{theorem}[Ticket one-more unforgeability]
After $q$ completed accepted blind-issuance sessions, producing more than $q$
distinct verifying ticket digests has probability at most the Blind-UOV
one-more-unforgeability advantage, plus the cross-message request-collision and
ticket-hash collision terms.
\end{theorem}

\begin{proof}
Map every completed issuance to the unique signed digest using
Lemma~\ref{lem:reqbind}.  More distinct verifying digests than completed
sessions then give more valid message/signature pairs than signing responses,
unless two payloads share a digest or Lemma~\ref{lem:buov-unique} fails.
\end{proof}

\section{Candidate concrete instantiations}

\subsection{Augmented well-built-request compiler for Blind-UOV}\label{sec:augcap}

The Blind-UOV framework already requires the user to prove that its masked
request is well built.  We strengthen that proved language instead of adding a
separate policy proof.  Let
\[
  R_{\mathsf{wb}}(y;m,r,\rho)=1
\]
denote Equations~\eqref{eq:b1}--\eqref{eq:b2}.  The native
well-built-request language is replaced by one augmented language:
\begin{align*}
 R_{\mathsf{aug}}(&\pp,\ctx,\mathsf{sid},\rid,y;\\
                  &M,r,\rho,k_{\mathrm{hold}},e)=1
 \quad\Longleftrightarrow\quad
 R_{\mathsf{issue}}(x;w)=1,
\end{align*}
where $\beta=y$ and Equation~\eqref{eq:i3} is $R_{\mathsf{wb}}$.  The issuer
executes the ordinary Blind-UOV response only
after verifying a post-quantum SE-NIZK argument for $R_{\mathsf{aug}}$.

\paragraph{Compiler interface.}
The integration replaces the native verification hook:
\[
 \mathsf{VerifyWellBuilt}(y,\pi_1)
 \quad\longrightarrow\quad
 \Piiss.\mathsf{Vfy}(\pp,\ctx,\mathsf{sid},\rid,y,
                    \pi_{\mathsf{issue}}).
\]
The response and finalization algorithms remain unmodified.  The spendable
object is one genuine Blind-UOV-III final signature.

\paragraph{Concrete circuit boundary.}
The implementation SHOULD expose the following fixed interface to the proof
compiler:
\begin{align*}
 \mathsf{pub}_{\mathsf{aug}}
   &= (\ctx,\mathsf{sid},\rid,\mathsf{tpk},\mathsf{bvk},y),\\
 \mathsf{wit}_{\mathsf{aug}}
   &= (\sn,k_{\mathrm{hold}},e,r,\rho).
\end{align*}
The circuit derives $C$, $h$, $M$, and $m$ internally rather than accepting
duplicate witness copies.  Its constraint system is the conjunction
\[
 \mathcal C_{\mathsf{aug}}=
 \mathcal C_{\mathsf{shape}}\wedge
 \mathcal C_{\mathsf{ticketHash}}\wedge
 \mathcal C_{\mathsf{BUOVmask}}\wedge
 \mathcal C_{\mathsf{holder}}\wedge
 \mathcal C_{\mathsf{trace}},
\]
where
\begin{align*}
 \mathcal C_{\mathsf{shape}} &: \mathsf{TicketShape}(M,\ctx)=1,\\
 \mathcal C_{\mathsf{ticketHash}} &: m=H_{\mathrm{ticket}}(\mathsf{Encode}(M)),\\
 \mathcal C_{\mathsf{BUOVmask}} &: c_r=\mathsf{CAP.Commit}(r;\rho)\ \wedge\
   y=r+H_{\BUOV}(m,c_r),\\
 \mathcal C_{\mathsf{holder}} &: h=H_{\mathrm{hold}}(k_{\mathrm{hold}}),\\
 \mathcal C_{\mathsf{trace}} &: R_{\mathsf{Ntr}}
   (\mathsf{tpk},C,\ctx,h;\rid,\sn,e)=1.
\end{align*}
The proof statement includes $\mathsf{sid}$ even though it is absent from the
ticket.  Consequently, replaying the same proof in a new issuance session
changes the statement and fails verification, while the final credential gains
no session link tag.

\begin{theorem}[Security inheritance of the augmented issuance proof]\label{thm:augcap}
Assume $\Piiss$ is post-quantum zero knowledge and simulation-extractable for
$R_{\mathsf{aug}}$.  Replacing the original well-built-request proof by
$\pi_{\mathsf{issue}}$ preserves Blind-UOV one-more unforgeability and
honest-protocol signer blindness, up to the SE-NIZK, cross-message request, and
hash-collision terms already shown in
Theorems~\ref{thm:certsound} and~\ref{thm:unlink}.
\end{theorem}

\begin{proof}
For one-more unforgeability, a forger against the augmented scheme is also a
forger against the original Blind-UOV scheme: the reduction verifies the
augmented proof itself and forwards each accepted $y$ to the original signer
instance.  If that interface requires the native well-built proof, the
reduction extracts $(m,r,\rho)$ and generates it.  Extraction failure is
charged to SE-AoK.  The augmented relation restricts the original well-built
language; the final signature format and verifier are unchanged.

For blindness, replace the augmented proof by a simulated proof.  The remaining
transcripts are ordinary Blind-UOV transcripts, so the native blindness game
applies directly.  The
public authenticated identity is session information but is absent from the
final ticket; threshold-encryption privacy hides it inside $M.C$.  Any attempt
to make the same masked request represent a different ticket digest triggers
Lemma~\ref{lem:buov-unique}, and any two payloads with one ticket digest trigger
a collision in $H_{\mathrm{ticket}}$.
\end{proof}

\paragraph{Corrected concretization in version 1.7.}
The request and witness are now fixed:
\begin{align*}
 \beta&=y,& b&=(r,\rho),\\
 m&=H_{\mathrm{ticket}}(\mathsf{Encode}(M)),&
 c_r&=\mathsf{CAP.Commit}(r;\rho),\\
 y&=r+H_{\BUOV}(m,c_r).&&
\end{align*}
The signer sees only the 72-byte value $y$, not $c_r$ or $m$.
Earlier drafts incorrectly exposed $c_r$ and claimed that CAP binding plus hash
collision resistance implied request uniqueness.  That argument fails when an
attacker chooses different commitments on both sides.  Version 1.6 restored the
paper ABI but then incorrectly multiplied the security of two independently
openable NIST-I lanes.  Version 1.7 instead uses a single 576-bit
Blind-UOV-III target and Theorem~\ref{thm:j-qrom}.  The removed
ticket commitment $c_M$ remains unnecessary: extraction supplies $M$, while
ticket-hash collision resistance prevents a second payload under one digest.

\paragraph{Executable refinement in version 1.8.}
Introduce the hidden 576-bit value
\[
 h_{\BUOV}=H_{\BUOV}(m,c_r).
\]
The executable incremental circuit now allocates $r$ and $h_{\BUOV}$ as secret
binary inputs, enforces their bitness, and materializes all 576 linear equations
$y_j+r_j+h_{\BUOV,j}=0$.  Its only remaining external assertion is the narrower
native subrelation
\[
 h_{\BUOV}=H_{\BUOV}(m,\mathsf{CAP.Commit}(r;\rho)),
\]
where $m$ is passed directly from the circuit-produced ticket-digest wires.
Thus a future native importer cannot choose a different $y$, $r$, hash image,
or ticket message without failing an ordinary circuit row.

\paragraph{Frozen Blind-UOV costing profile.}
Version 1.7 fixes the 31 October 2025 revision and its
\emph{Blind-UOV-III, Shorter, TCitH} row.  The paper gives security parameter
$\lambda=192$, $\bar n_r=576$ bits, $\bar n_x=1472$ bits, a 189.2~KB public
key, and an 11,644-byte final signature~\cite{blind-uov}.  The larger cached key
and signature buy a 576-bit request target whose generic quantum collision
scale exceeds the PQ-128 requirement.  Issuance remains the intentionally
expensive side of the trade-off.

\paragraph{Correct backend separation.}
Blind-UOV has two different proof objects.  Its CAP commitment and final CAP
proof are carried by the final blind signature.  Before the signer responds,
the user separately sends a NIZK $\pi_1$ proving that the masked request and
CAP commitment are well formed.  Our $\pi_{\mathsf{issue}}$ replaces and
augments this $\pi_1$; it does \emph{not} replace the final CAP proof.
Consequently the 11,644-byte final signature remains unchanged, while the
ticket/trace blocks are included only in the augmented issuance relation.

\paragraph{Native import contract and field boundary.}
For NIST-III/Shorter/TCitH, Table~2 of the paper fixes
$\tau=18$: two $2^{12}$-leaf trees and sixteen $2^{11}$-leaf trees,
$T_{\mathrm{open}}=174$, an explicit proof-of-work of 9 bits, and total
proof-of-work 13.9 bits.  The SNARK-friendly level-III profile uses an
eight-element Anemoi state over $\mathbb F_{2^{193}}$, with 240 R1CS
constraints per permutation~\cite{blind-uov}.  In particular, the abstract
$\rho$ contains the CAP salt and GGM-tree randomness; it is not the 32-byte
nonce used by the test adapter.

The paper specifies the algorithms but does not ship an executable CAP
constraint generator.  The new
\texttt{pq\_rbbc\_blind\_uov\_native.py} therefore defines a fail-closed import
contract.  Closure requires a source and parameter digest, a deterministic
native row-stream digest, zero external assertions, witness-independent
topology, verified Anemoi vectors, message/mask/hash-image wire identity, and
four independent tamper rejections.  It also requires the complete augmented
relation to be lifted from the portable $\mathbb F_2$ archive into
$\mathbb F_{2^{193}}$.  This structural certificate prevents accidental
under-constraint; it is not a substitute for cryptographic review of the CAP
implementation.  The official Anemoi repository supplies permutation and
equation-generation code, but not the paper's complete CAP importer
\cite{anemoi-code}.

For costing, use the paper's illustrative Shockwave/Anemoi NIST-I row only as
an anchor.  Blind-UOV-Is/Shorter is reported as approximately 16 million R1CS
constraints, a 2.8~MB proof, 32~s proving, and 2.5~s verification.  The paper
does not report the corresponding Blind-UOV-III row.  Scaling the native
constraint anchor by the witness-width ratio
$(576+1472)/(256+640)=2.2857$ gives a rough 36.57-million-row native baseline;
after the exact 684,419-row non-overlapping ticket/trace increment below, the
total is about 37.26 million.  The extra 1,152 bitness rows in the portable
v1.8 archive are already part of a genuine native $\pi_1$ relation and are not
double-counted in this combined estimate.
This is an extrapolation, not a paper result or local benchmark.  The paper
explicitly treats these as estimates and notes that the benchmarked
implementations do not provide the required zero-knowledge mode.  It also does
not establish the simulation-extractability required by our concurrent GCCA
proof.  Therefore Shockwave/Anemoi is a \emph{costing backend}, not yet the
security-frozen $\Piiss$ implementation~\cite{blind-uov}.

The multivariate commitment
$c=(\mu+F(r),G(r))$ and relation proofs of Bouillaguet et al.
are a promising later optimization because they are algebraically aligned with
UOV, but they add a newer MQ/bilinear-system assumption and are not the baseline
for the first auditable construction~\cite{mq-commitments}.

\subsection{Threshold Niederreiter for trace encryption}

\paragraph{Code and plaintext profile.}
Let $\mathcal G$ be a binary Goppa code of length $n$, dimension $k$, and
designed error weight $t$, with public systematic parity-check matrix
$H_{\mathrm{pub}}=(I_{n-k}\mid T)$.  The selected decoder MUST guarantee
$d_{\min}(\mathcal G)\geq 2t+1$.  Fix
$\rid\in\{0,1\}^{256}$ and $\sn\in\{0,1\}^{128}$, and define the fixed
48-byte plaintext
\[
  \mu=\mathsf{Encode}_{48}(\rid\|\sn).
\]
The online ciphertext uses a weight-$t$ error $e\in\mathbb F_2^n$ as its only
encryption witness.  Thus the abstract randomness $\rho_{\mathrm{tr}}$ is
instantiated by $e$.

\paragraph{Hash DEM.}
Sample $e$ uniformly from the weight-$t$ vectors and compute
\begin{align}
 s&=H_{\mathrm{pub}}e^\top\in\mathbb F_2^{n-k},
   &\mathrm{wt}(e)&=t, \tag{N1}\label{eq:n1}\\
 \mathsf{ad}&=\mathsf{Encode}(\ctx,\sn,h), &&\tag{N2}\label{eq:n2}\\
 Z&=\mathsf{SHAKE256}(
   \mathtt{PQ\mbox{-}RBBC/KDF}\|e\|s\|\ctx,640),
   &&\tag{N3}\label{eq:n3}\\
 (K_{\mathrm{mac}},P)&=Z,
   &|K_{\mathrm{mac}}|&=32,\quad |P|=48, \tag{N4}\label{eq:n4}\\
 c_\mu&=\mu\oplus P, &&\tag{N5}\label{eq:n5}\\
 \gamma&=\mathsf{KMAC256}_{K_{\mathrm{mac}}}
 (s\|c_\mu\|\mathsf{ad},256;
   \mathsf{custom}=\mathtt{PQ\mbox{-}RBBC/TAG}),
 &&\tag{N6}\label{eq:n6}\\
 C&=(s,c_\mu,\gamma). &&\tag{N7}\label{eq:n7}
\end{align}
The output length in (N3) is in bits.  SHAKE256 and KMAC256 are the standardized
extendable-output and MAC functions of FIPS~202 and SP~800-185,
respectively~\cite{fips202,sp800185}.  The composition $\HDEM$ is nevertheless a
new protocol profile, not a standardized AEAD and not yet a production claim.
A standardized one-time AEAD may replace it later without changing the abstract
$R_{\mathsf{enc}}$ interface.

\begin{definition}[Concrete trace relation]\label{def:ntr}
For $C=(s,c_\mu,\gamma)$, define
\[
 R_{\mathsf{Ntr}}(\mathsf{tpk},C,\ctx,h;
                    \rid,\sn,e)=1
\]
iff $\mathsf{tpk}$ selects $H_{\mathrm{pub}}$ and the fixed profile above, and
Equations~\eqref{eq:n1}--\eqref{eq:n7} all hold with
$\mu=\mathsf{Encode}_{48}(\rid\|\sn)$.  In this version,
\[
 \begin{aligned}
 &R_{\mathsf{enc}}(\mathsf{tpk},C,
     \mathsf{Encode}(\ctx,\sn,h);\rid\|\sn,\rho_{\mathrm{tr}})=1\\
 &\qquad\Longleftrightarrow\quad
 R_{\mathsf{Ntr}}(\mathsf{tpk},C,\ctx,h;\rid,\sn,e)=1,
 \qquad \rho_{\mathrm{tr}}=e.
 \end{aligned}
\]
Hence $\mathcal C_{\mathsf{trace}}$ is now the exact conjunction (N1)--(N7),
not an unspecified ``encrypt correctly'' box.
\end{definition}

\begin{lemma}[Unique syndrome witness]\label{lem:unique-error}
For a code with $d_{\min}\geq2t+1$, a syndrome $s$ has at most one witness
$e$ of weight at most $t$ satisfying $s=H_{\mathrm{pub}}e^\top$.
\end{lemma}

\begin{proof}
If $e$ and $e'$ were two such witnesses, then
$d=e+e'$ would satisfy $H_{\mathrm{pub}}d^\top=0$ and
$\mathrm{wt}(d)\leq2t$.  If $e\neq e'$, $d$ is a nonzero codeword below the
minimum distance, a contradiction.  Therefore $e=e'$.  This is the standard
unique-decoding argument used by Classic McEliece~\cite{classic-mceliece-spec}.
\end{proof}

\begin{corollary}[Unique trace opening]\label{cor:unique-trace}
The concrete relation $R_{\mathsf{Ntr}}$ has information-theoretically unique
plaintext openings.
\end{corollary}

\begin{proof}
For a fixed $C$, Lemma~\ref{lem:unique-error} fixes $e$ from $s$.  Equations
(N3)--(N4) then deterministically fix $P$ and $K_{\mathrm{mac}}$, and (N5)
fixes the sole $\mu=c_\mu\oplus P$.  The fixed-width encoding is injective, so
it fixes $(\rid,\sn)$.  Equation~(N6) additionally rejects any inconsistent
associated data.  No computational assumption is needed for uniqueness.
\end{proof}

\begin{proposition}[Concrete trace correctness]\label{prop:ntr-correct}
If threshold decoding reconstructs the unique error $e$ for $s$, the combiner
recomputes (N3)--(N6), returns $\mu=c_\mu\oplus P$, and accepts only if
$\gamma$ verifies and the decoded serial equals the visible $\sn$.
Consequently every honestly formed concrete trace ciphertext opens to exactly
$\rid\|\sn$.
\end{proposition}

\paragraph{Privacy assumption and proof boundary.}
Let $\mathsf{t\text{-}OW\text{-}SD}$ be the game in which an adversary receives
$H_{\mathrm{pub}}$, fewer than $\tau$ threshold shares, and
$s=H_{\mathrm{pub}}e^\top$ for uniform weight-$t$ $e$, and must output $e$.
We assume this advantage is negligible for the frozen distributed decoder.

\begin{theorem}[Conditional concrete trace privacy]\label{thm:ntr-privacy}
Model SHAKE256 as a quantum random oracle, assume KMAC256 is a quantum-secure
PRF, never reuse $e$, and let $\epsilon_{\mathsf{O2H}}(q_H)$ denote the loss of
the applicable one-way-to-hiding lemma for $q_H$ quantum hash queries.  Then for
every fewer-than-threshold $\QPT$ distinguisher $\mathcal A$ there are $\QPT$
reductions $\mathcal B_1,\mathcal B_2$ such that
\begin{align*}
 \Adv^{\mathsf{t\text{-}IND\text{-}CPA}}_{\mathsf{Ntr},\mathcal A}
 \leq{}&
 \Adv^{\mathsf{t\text{-}OW\text{-}SD}}_{\mathcal G,\mathcal B_1}
 +\epsilon_{\mathsf{O2H}}(q_H)
 +\Adv^{\mathsf{qPRF}}_{\mathsf{KMAC256},\mathcal B_2}.
\end{align*}
\end{theorem}

\begin{proof}
Starting from the real game, replace the 640-bit SHAKE256 value at the hidden
input $(e,s,\ctx)$ by a uniform string.  A distinguisher for this step
either yields a threshold syndrome-decoding solver or is bounded by the stated
one-way-to-hiding loss.  Now $P$ is an independent uniform 48-byte pad and
$K_{\mathrm{mac}}$ is an independent uniform key.  Replace KMAC256 under that
key by a random function, paying the qPRF term.  The value
$c_\mu=\mu_b\oplus P$ and the tag are then independent of $b$, so the final
advantage is zero.  Summing the hybrid differences gives the bound.
\end{proof}

A new accepted tuple $(s,c'_\mu,\gamma')$ with the same hidden derived key and
different authenticated input gives a KMAC forgery, up to the same
syndrome-decoding/QROM bad event.  Thus the xor pad gives one-time privacy and
KMAC256 gives one-time ciphertext integrity.  Substituting
Theorem~\ref{thm:ntr-privacy} for the abstract $\tau$-IND-CPA term in
Theorem~\ref{thm:gcca} yields protocol-level $\tau$-GCCA, plus the SE-AoK,
Blind-UOV, CAP, unique-opening, authorization, and hash terms already displayed
there.  This does \emph{not} prove standard IND-CCA security for a bare
ciphertext endpoint.

The threshold-Niederreiter work solves malformed-ciphertext attacks by attaching
a straight-line-extractable validity proof to every public ciphertext; its
authors report about 30.5~KB for a three-party ciphertext, with the proof
dominating the syndrome~\cite{threshold-niederreiter,threshold-niederreiter-slides}.
Here the same validity language
$\{s:\exists e,\mathrm{wt}(e)=t,\;s=He^\top\}$
is certified once inside $\pi_{\mathsf{issue}}$.  The online ticket retains only
$C=(s,c_\mu,\gamma)$, and the raw distributed decoder remains unreachable.

\paragraph{Exact online byte counts.}
Because $|c_\mu|=48$ and $|\gamma|=32$, the trace ciphertext is exactly the
syndrome length plus 80 bytes.  Using the published Classic McEliece
implementation sizes gives the following anchors~\cite{classic-mceliece-impl}:
\begin{center}
\begin{tabular}{lrrr}
\toprule
Parameter anchor & Public key & Syndrome & This $C_{\mathrm{trace}}$\\
\midrule
mceliece348864  & 261120 B  &  96 B & 176 B\\
mceliece460896  & 524160 B  & 156 B & 236 B\\
mceliece6688128 & 1044992 B & 208 B & 288 B\\
mceliece6960119 & 1047319 B & 194 B & 274 B\\
mceliece8192128 & 1357824 B & 208 B & 288 B\\
\bottomrule
\end{tabular}
\end{center}
The aggressive working anchor is mceliece6688128: approximately 1.045~MB of
cached public key buys a 288-byte trace ciphertext.  The Blind-UOV-III
profile and exact 12,012-byte complete ticket are derived in
Section~\ref{sec:circuit-audit}; these are still research targets rather than
implementation benchmarks.

\begin{remark}[Mandatory 2026 cryptanalytic hold]\label{rem:mceliece-hold}
The anchor above is not frozen.  A 10 August 2026 preprint gives a classical
quasipolynomial-time distinguisher applying to all Classic McEliece NIST
parameter sets and extends it to a heuristic quasipolynomial-time decryption
attack; its concrete improvements are described as not yet practical
\cite{mceliece-2026-attack}.  That is not a demonstrated practical break, but it
materially changes confidence.  Deployment parameter selection is blocked
until independent review quantifies this attack for the threshold construction.
\end{remark}

\subsection{Reproducible five-block circuit audit}\label{sec:circuit-audit}

\paragraph{Counting model.}
The paper's Anemoi profile uses an R1CS over $\mathbb F_{2^{129}}$ for NIST
level I.  Addition, xor, multiplication by public constants, fixed-width
concatenation, and the public binary syndrome map are linear.  Each Boolean AND
or field multiplication contributes one nonlinear R1CS constraint.  Following
the Blind-UOV paper, one Keccak-f permutation costs
$24\cdot1600=38{,}400$ constraints~\cite{blind-uov}.  The audit counts only
constraints added beyond the native Blind-UOV-III well-built-request $\pi_1$
relation.  Because the paper reports an R1CS anchor only for NIST-I, the
level-III baseline below is explicitly extrapolated.

\paragraph{Exact-weight gadget.}
The witness $e$ has 6,688 bits.  Pad it with 1,504 public zeros to 8,192 inputs
and sum them using a fixed balanced binary adder tree.  Combining two $j$-bit
counters uses one half-adder and $j-1$ full adders.  Since xor is linear, the
nonlinear cost is $1+2(j-1)=2j-1$.  The exact cost is therefore
\[
 \sum_{j=1}^{13}\frac{8192}{2^j}(2j-1)=24{,}547
\]
AND constraints.  Equating the final 14-bit count to 128 is linear.  This
gadget avoids the invalid shortcut $\sum e_i=128$ in characteristic two, which
would check only parity.  Separately, $e$, $\sn$, and $k_{\mathrm{hold}}$
require $6688+128+256=7{,}072$ bitness constraints.

\paragraph{Exact Keccak accounting.}
SHAKE256 has a 136-byte rate.  With pad10*1, an $L$-byte message uses
$\lfloor L/136\rfloor+1$ absorb permutations.  The fixed encodings give:
\begin{center}
\begin{tabular}{lrrr}
\toprule
Invocation & Absorbed bytes & Keccak-f calls & Nonlinear constraints\\
\midrule
$H_{\mathrm{ticket}}(M)$ & $14+368=382$ & 3 & 115,200\\
$H_{\mathrm{hold}}(k)$ & $12+32=44$ & 1 & 38,400\\
Trace SHAKE KDF & $11+836+208+32=1087$ & 8 & 307,200\\
KMAC256 & $136+136+336+3=611$ & 5 & 192,000\\
\midrule
Total & & 17 & 652,800\\
\bottomrule
\end{tabular}
\end{center}
For KMAC, the first 136-byte block is the cSHAKE function/customization prefix,
the second is $\mathsf{bytepad}(\mathsf{encode\_string}(K),136)$, the
authenticated data is $208+48+80=336$ bytes, and
$\mathsf{right\_encode}(256)$ is three bytes.  All squeezing fits in the first
rate block and adds no permutation.

\paragraph{Five-block result.}
The incremental nonlinear counts are:
\begin{center}
\begin{tabular}{lrl}
\toprule
Block & Constraints & Contents\\
\midrule
$\mathcal C_{\mathsf{shape}}$ & 128 & serial bitness; widths/wiring linear\\
$\mathcal C_{\mathsf{ticketHash}}$ & 115,200 & three SHAKE permutations\\
$\mathcal C_{\mathsf{BUOVmask}}$ & 1,152 & bitness of $r,h_{\BUOV}$; 576 mask equations linear\\
$\mathcal C_{\mathsf{holder}}$ & 38,656 & 256 bitness + one SHAKE permutation\\
$\mathcal C_{\mathsf{trace}}$ & 530,435 & $e$ bitness, weight, KDF, and KMAC\\
\midrule
Portable v1.8 relation & 685,571 & includes 1,152 native-overlap hardening rows\\
\bottomrule
\end{tabular}
\end{center}
Of the portable total, 1,152 bitness rows protect the explicit $r$ and
$h_{\BUOV}$ interface while the native generator is absent.  A genuine native
$\pi_1$ already constrains those binary values, so only 684,419 rows are
non-overlapping additions to the paper-anchored native baseline.  The 576
$y=r+h_{\BUOV}$ equations are linear.
The syndrome equation adds 1,664 linear output equalities but no nonlinear
gate.  The xor pad, associated-data construction, external/internal serial
equality, and final constant comparisons are also linear.

\paragraph{Paper-anchored issuance estimate.}
The level-III native baseline is estimated from the paper's rounded
16-million-constraint NIST-I figure by the witness-width ratio:
\[
 N_{\mathsf{native,III}}\approx16{,}000{,}000
 \frac{576+1472}{256+640}=36{,}571{,}429.
\]
Adding the exact incremental relation gives
\[
 N_{\pi_{\mathsf{issue}}}\approx36{,}571{,}429+684{,}419
 =37{,}255{,}848,
\]
or 2.3285 times the NIST-I anchor.  Applying the paper's square-root
proof-size scaling and linear time scaling gives a rough 4.27~MB proof,
74.5~s proving, and 5.82~s verification.  None of these level-III figures is a
paper result or local benchmark; the scaling may omit parameter-dependent
costs.  Peak prover memory remains a measurement obligation.

\begin{proposition}[Monolithic-circuit feasibility decision]
Under the costing model above, adding the entire ticket and trace relation to
the estimated Blind-UOV-III well-built-request NIZK changes its baseline by
about 1.87\%.
There is therefore no communication-based reason to split
$R_{\mathsf{aug}}$ into linked proofs at this stage.  A split is justified only
if the selected simulation-extractable backend, implementation memory, or
prover parallelism makes the monolithic circuit impractical.
\end{proposition}

\paragraph{Exact online object under the frozen profile.}
For mceliece6688128, $|C|=288$ bytes and
\[
 |M|=32+16+32+288=368\text{ bytes}.
\]
Combining this with the published 11,644-byte Blind-UOV-III/Shorter signature
gives
\[
 |T|=|M|+|\sigma|=12{,}012\text{ bytes}
\]
before transport framing.  The cached public material is approximately
1,044,992 bytes for the trace key plus the paper's rounded 189.2~KB Blind-UOV
verification key.  Relative to a 70~KB online credential,
the frozen target is about 5.83 times smaller.  The roughly 4.27~MB
issuance NIZK is offline and
is never forwarded to a satellite verifier.

\begin{remark}[Audit reproducibility]
The companion executable relation recomputes every exact incremental integer
above and records it in the version-1.7 manifest.  Native level-III estimates
are kept in separately labelled fields so a rounded extrapolation cannot be
mistaken for a benchmark.
\end{remark}

\subsection{Executable reference relation and rejection tests}
\label{sec:executable-relation}

Version 1.8 extends the incremental relation in
\texttt{pq\_rbbc\_reference.py}.  The witness builder derives
$h$, $s$, $c_\mu$, $\gamma$, the canonical 368-byte payload, and its ticket
digest from $(\rid,\ctx,\sn,k_{\mathrm{hold}},e)$.  The circuit builder then
emits a stream of characteristic-two events of four kinds: linear wire
definitions, linear equality assertions, Boolean multiplication gates, and
bitness constraints $x(x+1)=0$.  A backend can consume the same stream without
retaining the whole circuit in memory.  This version executes the stream and
counts it, but does not yet flatten it into a proof system's matrix format.

The symbolic Keccak-f implementation performs all 24 rounds.  Its SHAKE256
output is checked against Python's independent FIPS~202 implementation, and
the SP~800-185 KMAC256 output is checked against the OpenSSL 3 KMAC provider
\cite{fips202,sp800185}.
For the frozen honest vector the generated circuit reports:
\begin{center}
\begin{tabular}{lr}
\toprule
Measured generator quantity & Value\\
\midrule
Allocated wire identifiers & 2,980,304\\
Linear definitions & 2,282,475\\
Linear equality assertions & 3,534\\
Nonlinear constraints & 685,571\\
Public input bits & 4,032\\
Secret input bits & 8,224\\
External Blind-UOV assertions & 1\\
Keccak-f permutations & 17\\
Failed assertions & 0\\
\bottomrule
\end{tabular}
\end{center}
Thus the executable result exactly reproduces every nonlinear block total in
the preceding audit.  The linear-definition figure is diagnostic IR size, not
an additional nonlinear count: the costing model requires a backend to inline
affine expressions.  A backend that instead materializes each affine wire as
an R1CS row must report that larger total explicitly.  On the present reference
environment, one honest streaming generation takes about six seconds.  This is Python relation
generation time, not NIZK proving time, and no prover-memory claim follows from
it.

The regression suite accepts the honest witness and sends eleven independent
mutations through the complete circuit: wrong error weight, changed syndrome,
changed masked identity, changed holder hash, changed KMAC tag, changed serial,
changed blind request, changed blind mask, changed Blind-UOV hash image, changed
CAP randomness, and changed common context.  All eleven are rejected.
The checked output is recorded in
\path{pq_rbbc_reference_manifest_v1_8.json}.

\paragraph{Exact implementation boundary.}
The $H=(I\mid T)$ matrix in the reference vector is generated deterministically
from a public seed and is only a rank-controlled test fixture; it is not a
certified binary Goppa key.  Similarly,
\texttt{TestBlindUOVAdapter} checks the CAP-hash subrelation with a deterministic
SHAKE construction, but is not a blind signature.  It marks one native
Blind-UOV-III CAP.Commit-plus-hash boundary whose nonlinear cost remains in the
paper-anchored baseline.  The 576 public $y$ bits and 1,152 secret
$r,h_{\BUOV}$ bits are serialized as circuit inputs; $m$ is derived internally,
and $\rho,c_r$ remain hidden behind the native boundary.  Consequently this implementation establishes that the incremental
relation is internally consistent and executable; it does not establish that
the combined native Blind-UOV proof is zero knowledge or simulation
extractable.

\subsection{Portable binary R1CS lowering}\label{sec:br1cs}

Version 1.8 uses the proof-backend exchange layer in
\texttt{pq\_rbbc\_br1cs.py}.  It consumes the same event stream as the counting
sink and serializes three standard rank-one forms over $\mathbb F_2$:
\begin{align*}
  \mathsf{MUL}(a,b,c)&:\quad a\cdot b=c,\\
  \mathsf{LINDEF}(c;x_1,\ldots,x_j;k)&:\quad
       (c+x_1+\cdots+x_j+k)\cdot 1=0,\\
  \mathsf{ASSERT}(x_1,\ldots,x_j;k)&:\quad
       (x_1+\cdots+x_j+k)\cdot 1=0.
\end{align*}
Wire identifiers and arities use unsigned LEB128 encodings.  A fixed 124-byte
header commits to all row counts, the body length, and the SHA-256 digest of the
body.  Public and secret input declarations contain identifiers and visibility
only; assignment values are not embedded in the circuit archive.

The resulting honest archive has the following exact dimensions:
\begin{center}
\begin{tabular}{lr}
\toprule
Portable backend quantity & Value\\
\midrule
Wires & 2,980,304\\
Public input bits & 4,032\\
Secret input bits & 8,224\\
Nonlinear R1CS rows & 685,571\\
Materialized affine-definition rows & 2,282,475\\
Linear assertion rows & 3,534\\
Total R1CS rows & 2,971,580\\
Archive size & 49,227,687 bytes\\
\bottomrule
\end{tabular}
\end{center}
The 685,571 figure is therefore a \emph{nonlinear-only} cost under the affine
elimination model used in Section~\ref{sec:circuit-audit}.  The 2,971,580 figure
is the portable total when every affine definition is materialized as its own
rank-one row.  They are not competing measurements.  A concrete proof backend
may recover the smaller representation by substituting affine expressions into
multiplication rows, but that optimization requires its own soundness argument
and resource benchmark.

Although the exchange format names $\mathbb F_2$ as its base relation, the same
Boolean assignment and equations embed into $\mathbb F_{2^d}$: addition remains
xor, and the input bitness equations prevent non-Boolean witnesses.  This makes
the archive compatible in principle with both the $\mathbb F_{2^{129}}$
costing anchor and the target $\mathbb F_{2^{193}}$ level-III field, without
claiming that the required lift or a specific prover has been integrated.

An independent reader reconstructs and evaluates all 2,971,580 rows.  It
accepts the honest assignment, rejects a one-bit assignment mutation, and
rejects a one-bit archive corruption through the body digest.  Lowering a
wrong-weight witness produces exactly the same 49,227,563-byte body and the
same digest as the honest witness, showing that circuit structure is
witness-independent; its assignment nevertheless fails 320 rows.  Five
end-to-end backend tests pass; wall time is environment-dependent.  The checked
metadata is recorded in
\texttt{pq\_rbbc\_br1cs\_manifest\_v1\_8.json}.

\paragraph{Remaining external rows.}
The archive records one named external assertion for
$h_{\BUOV}=H_{\BUOV}(m,\mathsf{CAP.Commit}(r;\rho))$, but cannot evaluate it.
The linear equation $y=r+h_{\BUOV}$ is no longer external.  This narrowed
boundary is deliberate and machine-visible: an archive with an unchecked
external assertion is not a closed issuance proof.  The next integration must
replace that marker with the actual TCitH/Anemoi constraints; deleting or
treating it as automatically true is forbidden.

\subsection{Local mixing as a future trace-encryption candidate}

The core does not depend on elliptic curves, lattices, or codes.  A future
local-mixing encryption may replace threshold Niederreiter only if it supplies
all four properties (T1)--(T4), an efficiently provable NP encryption relation,
and a practical $\tau$-party opening protocol.  A small standalone ciphertext
claim is insufficient by itself: threshold opening, unique opening, and
extractable issuance binding are the missing system properties.

\section{What has and has not been proved}

\begin{center}
\begin{tabular}{p{0.31\linewidth}p{0.61\linewidth}}
\toprule
\textbf{Status} & \textbf{Result or obligation}\\
\midrule
Conditional abstract reductions & Correctness; proof/request/message binding;
certified ciphertext validity; signature-gated decoder safety; issuer
unlinkability against an honest-protocol curious issuer; trace soundness;
one-more ticket unforgeability; and $\tau$-GCCA trace privacy.\\
\addlinespace
Assumed component property & Post-quantum SE-NIZK AoK; CAP straightline
extraction and unique committed mask; QROM cross-message request binding;
Blind-UOV blindness and one-more
unforgeability; threshold syndrome-decoding privacy with robust shares; QROM
one-way-to-hiding; KMAC integrity; and
authorization unforgeability.\\
\addlinespace
Concretely instantiated & Blind-UOV-III/Shorter size profile; executable
five-block augmented relation; exact Niederreiter--SHAKE--KMAC trace relation;
information-theoretic unique opening; independently checked SHAKE/KMAC;
full-circuit positive and negative tests; an internally constrained
$y=r+h_{\BUOV}$ equation; a fail-closed native-import contract; a checksummed
portable R1CS archive with independent row evaluation; 685,571 incremental
nonlinear rows; 2,971,580 materialized total rows; a 72-byte issuance request
excluding proof; and a 12,012-byte online-ticket target.\\
\addlinespace
Not yet instantiated & A post-quantum zero-knowledge and
simulation-extractable implementation of $\Piiss$; a native
$\mathbb F_{2^{193}}$ Blind-UOV-III CAP-hash circuit import and validation of
the CAP properties used by Theorem~\ref{thm:j-qrom};
proof-system-specific affine elimination; import of a real Goppa
key; peak prover memory; robust threshold-share transcript; proof benchmarks;
and a reviewed response to the August 2026 McEliece cryptanalysis.\\
\addlinespace
Explicitly not claimed & Standard IND-CCA security of the bare ciphertext,
security against a malicious issuer, metadata/timing anonymity without system
controls, UC composability, side-channel resistance, satellite AKE security,
or FLEO/handover security.\\
\bottomrule
\end{tabular}
\end{center}

\section{Implementation invariants}

The following invariants are part of the cryptographic design and should be
machine-tested or enforced by type/API boundaries.

\begin{enumerate}[label=INV-\arabic*,leftmargin=*]
  \item There is no production function that accepts a bare trace ciphertext
  and returns a partial decryption.
  \item The bytes hashed for the issuer signature are exactly the canonical
  bytes parsed by ticket verification and opening.
  \item The issuer signs only after successful proof verification and records a
  completed session only after returning its blind response.
  \item The blind request proven in $\Piiss$ is byte-for-byte the request passed
  to the blind-signing implementation.
  \item The public issuance request contains exactly the 72-byte value $y$; it
  contains neither $m$ nor $c_r$.  The mask and commitment randomness remain
  hidden and the commit-before-hash order is enforced.
  \item The opening authorization includes the ticket digest and evidence
  digest, and replay is statefully rejected.
  \item The reconstructed serial number must equal the clear serial number.
  \item Issuer keys and opening keys have independent generation ceremonies,
  storage, rotation, and compromise domains.
  \item Epoch metadata is federation-fixed; the issuer cannot personalize it.
  \item Every component and parameter set has an explicit post-quantum security
  claim at the chosen target category.
  \item The parser accepts exactly the frozen 368-byte payload layout and
  rejects variable-length or alternate encodings before signature verification.
  \item The issuance circuit enforces bitness of all 6,688 error coordinates
  and proves exact weight 128 with the specified counter gadget; a parity-only
  check is forbidden.
  \item The Shockwave/Anemoi figures may be used for costing only until a
  backend meeting the stated zero-knowledge and simulation-extractability model
  is proved or substituted.
  \item Production configuration rejects the deterministic parity-check matrix
  and test-only blind-signature adapter; neither fixture may cross the backend
  boundary.
  \item Backend lowering must preserve every multiplication and equality event,
  eliminate affine wires soundly, and accept/reject the published reference and
  mutation vectors identically to the executable IR.
  \item A serialized circuit is accepted only if its declared length, row
  counts, sequential wire discipline, and body digest all verify before proof
  generation or verification.
  \item Any nonzero external-assertion count makes the archive incomplete for
  production issuance; the Blind-UOV-III marker must be replaced, never ignored.
  \item The 576 mask bits and 576 Blind-UOV hash-image bits are secret circuit
  inputs with bitness constraints; all 576 equations $y_j+r_j+h_{\BUOV,j}=0$
  must remain ordinary rows after lowering.
  \item A native import is accepted only if it matches the frozen
  $\mathbb F_{2^{193}}$ profile, has witness-independent topology and zero
  external assertions, and satisfies every check in the native-import contract.
\end{enumerate}

\section{Next proof step}

The incremental relation, linear mask binding, unique opening, fixed layout,
circuit audit, executable witness builder, rejection tests, portable R1CS
lowering, and native-import contract are now complete.  The next step has two
coupled proof obligations:

\begin{quote}
\emph{Implement or obtain a bit-exact TCitH/Anemoi generator for Protocols
8--10, pin its source and parameter files, and compile its CAP.Commit-plus-hash
subrelation over $\mathbb F_{2^{193}}$.  Feed it the exact digest, mask, and
hash-image wires already named by the import contract; the resulting combined
archive must contain zero external assertions and pass all four native tamper
tests.  Validate the straightline-extraction, unique-mask, commit-before-hash,
and domain-separation conditions used in Theorem~\ref{thm:j-qrom}.  Then select a concrete backend,
prove any affine-elimination pass sound, and establish post-quantum zero
knowledge and simulation extractability in the model required by the GCCA
reduction.}
\end{quote}

The current count favors one monolithic proof; splitting is no longer the
default.  The Niederreiter deployment parameter MUST remain provisional under
Remark~\ref{rem:mceliece-hold}.  After backend qualification, the remaining
cryptographic component is the robust threshold-share transcript.  Satellite
AKE integration remains out of scope until those two obligations close.

\begin{thebibliography}{99}

\bibitem{blind-uov}
Charles Bouillaguet, Thibauld Feneuil, Jules Maire, Matthieu Rivain,
Julia Sauvage, and Damien Vergnaud,
\emph{Blinding Post-Quantum Hash-and-Sign Signatures},
IACR Cryptology ePrint Archive, Report 2025/895, 2025.
\url{https://eprint.iacr.org/2025/895}

\bibitem{anemoi-code}
Cl\'emence Bouvier et al.,
\emph{Anemoi: a Family of ZK-friendly Arithmetization-Oriented Hash
Functions}, reference Sage implementation.
\url{https://github.com/anemoi-hash/anemoi-hash}

\bibitem{bht}
Gilles Brassard, Peter H\o yer, and Alain Tapp,
\emph{Quantum Algorithm for the Collision Problem},
arXiv:quant-ph/9705002, 1997.
\url{https://arxiv.org/abs/quant-ph/9705002}

\bibitem{compressed-oracle}
Kai-Min Chung, Serge Fehr, Yu-Hsuan Huang, and Tai-Ning Liao,
\emph{On the Compressed-Oracle Technique, and Post-Quantum Security of Proofs
of Sequential Work},
Advances in Cryptology---EUROCRYPT 2021; arXiv:2010.11658.
\url{https://arxiv.org/abs/2010.11658}

\bibitem{blind-uov-slides}
Thibauld Feneuil,
\emph{Blinding Post-Quantum Hash-and-Sign Signatures},
IEEE Symposium on Security and Privacy presentation slides, 2026.
\url{https://www.thibauld-feneuil.fr/docs/talks/2026-05-19_SP26_blind-sigs.pdf}

\bibitem{mq-commitments}
Charles Bouillaguet, Thibauld Feneuil, Jules Maire, Matthieu Rivain,
Julia Sauvage, and Damien Vergnaud,
\emph{Multivariate Commitments and Signatures with Efficient Protocols},
IACR Cryptology ePrint Archive, Report 2025/2035, 2025.
\url{https://eprint.iacr.org/2025/2035}

\bibitem{threshold-niederreiter}
Pascal Giorgi, Fabien Laguillaumie, Lucas Ottow, and Damien Vergnaud,
\emph{Threshold Niederreiter: Chosen-Ciphertext Security and Improved
Distributed Decoding},
IACR Cryptology ePrint Archive, Report 2025/757, 2025.
\url{https://eprint.iacr.org/2025/757}

\bibitem{threshold-niederreiter-slides}
Lucas Ottow,
\emph{Threshold Niederreiter: Chosen-Ciphertext Security and Improved
Distributed Decoding}, presentation slides, 2025.
\url{https://indico.math.cnrs.fr/event/11948/attachments/5546/9681/Lucas_Ottow.pdf}

\bibitem{classic-mceliece-spec}
Daniel J. Bernstein et al.,
\emph{Classic McEliece: conservative code-based cryptography}, specification.
\url{https://classic.mceliece.org/mceliece-spec-20221023.pdf}

\bibitem{classic-mceliece-impl}
Classic McEliece team,
\emph{Implementations and parameter sizes}.
\url{https://classic.mceliece.org/impl.html}

\bibitem{fips202}
National Institute of Standards and Technology,
\emph{FIPS 202: SHA-3 Standard}, 2015.
\url{https://csrc.nist.gov/pubs/fips/202/final}

\bibitem{sp800185}
National Institute of Standards and Technology,
\emph{SP 800-185: SHA-3 Derived Functions}, 2016.
\url{https://csrc.nist.gov/pubs/sp/800/185/final}

\bibitem{mceliece-2026-attack}
Ashrujit Ghoshal, Yuval Ishai, Aayush Jain, and Nuozhou Sun,
\emph{Quasipolynomial Cryptanalysis of the McEliece Cryptosystem
(or: PIR Meets McEliece)}, IACR Cryptology ePrint Archive,
Report 2026/1630, 2026.
\url{https://eprint.iacr.org/2026/1630}

\bibitem{blind-ciphertext-signatures}
Olivier Blazy, Georg Fuchsbauer, David Pointcheval, and Damien Vergnaud,
\emph{Signatures on Randomizable Ciphertexts},
Public Key Cryptography---PKC 2011.
\url{https://www.blazy.eu/Papers/2011_pkc.pdf}

\end{thebibliography}

\end{document}
